
\section{Introduction to Github}

GitHub is a Git repository web-based hosting service which offers all of the functionality of Git as well as adding many of its own features. Unlike Git which is strictly a command-line tool, Github provides a web-based graphical interface and desktop as well as mobile integration. It also provides access control and several collaboration features such as wikis, task management, and bug tracking and feature requests for every project.\\

\begin{figure}[H]
	\centering
	\includegraphics[width=0.7\linewidth]{images/Github}
	\caption{Github Logo}
	\label{fig:930px-githublogo2013padded}
\end{figure}


This means that you can do things like:
\begin{itemize}
\item Frictionless Context Switching.\\ Create a branch to try out an
idea, commit a few times, switch back to where you branched from,
apply a patch, switch back to where you are experimenting, and merge
it in.
\item Role-Based Code lines. \\ Have a branch that always contains only
what goes to production, another that you merge work into for testing,
and several smaller ones for day to day work.
\item Feature Based Work flow. \\ Create new branches for each new
feature you're working on so you can seamlessly switch back and forth
between them, then delete each branch when that feature gets merged
into your main line.
\item Disposable Experimentation.\\  Create a branch to experiment in,
realize it's not going to work, and just delete it - abandoning the
work with nobody else ever seeing it (even if you've pushed other
branches in the meantime).
\end{itemize}
Notably, when you push to a remote repository, you do not have to push
all of your branches. You can choose to share just one of your
branches, a few of them, or all of them. This tends to free people to
try new ideas without worrying about having to plan how and when they
are going to merge it in or share it with others.\\ \\

\subsection{What is Git?}
\begin{figure}[H]
\centering
\includegraphics[width=0.3\textwidth]{images/git.png}                   
\caption{Git Logo}
\hspace{-1.5em}
\end{figure}
Git is a distributed revision control and source code management (SCM) system with an emphasis on speed, data integrity, and support for distributed, non-linear workflows. Git was initially designed and developed by Linus Torvalds for Linux kernel development in 2005, and has since become the most widely adopted version control system for software development.\\

As with most other distributed revision control systems, and unlike most client–server systems, every Git working directory is a full-fledged repository with complete history and full version-tracking capabilities, independent of network access or a central server. Like the Linux kernel, Git is free and open source software distributed under the terms of the GNU General Public License version 2 to handle everything from small to very large projects with speed and efficiency.\\
